% Results chapter. Two completed forecasting ablations, each a pre-registered comparison
% adjudicated by the model confidence set of the methodology chapter: the S5 multi-venue
% group in-context learning study (report 47) and the S6 lead-matched weather study (report
% 48). Written to be \input into the thesis body; it reuses the bibliography keys of the
% methodology chapter (chapters/ref.bib), in particular hansen_model_2011.

\chapter{Results}
\label{chap:results}

Two questions about the flagship forecaster admit a clean experimental answer, and this
chapter reports both. The first is whether forecasting the estate's venues together, so that
a foundation model can borrow strength across them, beats forecasting each alone. The second
is how much of the exogenous weather path's apparent value is real once the weather is placed
on the footing the model would actually have at serving. Both are pre-registered comparisons,
both are adjudicated by the model confidence set of Chapter~\ref{chap:methodology}, and both
return a negative result at full prominence rather than a tuned win. Reporting them together
is deliberate: they share an instrument and a discipline, and each is a case where the honest
finding is that an appealing elaboration of the forecaster buys nothing the evidence can
distinguish.

\section{Cross-series in-context learning does not beat the univariate baseline}
\label{sec:group-icl}

Chronos-2 supports cross-series learning: several related series passed together in one call
condition on one another, so a data-rich venue can inform a data-poor one. The serving path
did not use this, passing every venue under a single shared identifier one call at a time, so
the venues were forecast as isolated univariate series. The hypothesis, recorded before the
run, was that grouping would help most at the sparsest venue, the village hall at Ellel, which
has the most to gain from borrowing a richer series, and little at the Beer Hall, whose own
history is long. Three arms were compared on the same rolling origins: the univariate baseline,
the Beer Hall and Ellel grouped, and those two joined by Two River Taps carrying constructed
post-closure zeros.

The load-bearing precaution is a leakage guard, and it is worth stating why it is not a
convention but a hard check. Under cross-series attention, a boundary error in one series
contaminates the others: if any venue in a grouped context carried a value dated at or after
the forecast origin, the model would be reading the future of the whole panel, and the
resulting inflation would be invisible in the summary statistics. The guard therefore truncates
every series in every grouped context at the origin, the target venue and the non-target venues
alike, and raises rather than warns if any context reaches beyond it. A second property is
enforced rather than assumed: because the pipeline mixes exactly the series that share a
processing batch, the cross-learning group is made to coincide with that batch by pinning its
size to the number of venues and ordering the frame so each origin forms its own isolated group.
Batching several origins into one call is then numerically identical to issuing them separately,
which was verified on the model to a maximum absolute difference of zero, so throughput never
changes a forecast.

The result is that grouping does not help anywhere the evidence can distinguish, and where it
moves the number it moves it the wrong way. On the two hundred and sixty aligned origins the
grouped arms are indistinguishable from or worse than the univariate baseline. At the Beer Hall
the univariate arm has the lowest mean scaled error and grouping is a small but statistically
significant loss, the paired bootstrap interval on the difference excluding zero, though all
three arms sit inside the ninety percent confidence set. At Ellel the fuller group is
directionally the best, in the predicted order, but the improvement is not significant, every
paired interval spanning zero. At Two River Taps the grouped arm is eliminated from the
confidence set outright. The pre-registered adoption rule required a grouped arm both to enter
the ninety percent set and to have the lower mean at the venue concerned; no arm satisfies both
at a venue where the criterion applies, so nothing is adopted and no served model changes. The
capability is real, a grouped forecast differs materially from the univariate one, but on this
estate it is not an improvement.

\section{Perfect weather is no better than a seven-day-ahead forecast of it}
\label{sec:weather}

The Beer Hall's served model conditions on weather, and every accuracy figure it has produced
was measured against weather archived close to the valid time. At serving, a forecast made at
one origin needs weather for the seven days that follow, so the seventh day is a seven-day-ahead
weather forecast with seven-day-ahead skill. Evaluating the model against weather archived one
day out therefore hands it better covariates than it will ever have in production. This is not
target leakage, and the distinction matters enough to name precisely: it is a mismatch in
covariate quality between evaluation and serving, and it biases the exogenous path's apparent
value upward. The study asks how much of that value survives when the weather is placed on the
basis the model would actually have.

Five arms were compared, differing only in the seven-day window of weather handed to the model
and holding the training-context weather and every non-weather covariate fixed: no weather at
all; observed reanalysis at the valid time, an unachievable upper bound of perfect foreknowledge;
the archived-near-valid-time basis that every published figure used; a fixed operational lead;
and the horizon-matched basis in which the value for horizon step $h$ is the forecast issued $h$
days earlier, which is what serving delivers. The lead-matched arms were built from a single
pinned global weather model, because the automatic selection populates the longer leads from a
local high-resolution model that leaves them null, and a basis built on partly-null leads would
measure missingness rather than skill. A gate analogous to the panel-leakage guard of the
previous section enforces that no arm can supply a weather value whose lead is shorter than its
horizon step, verified by value against the pinned store and proven to fire on a planted
short-lead.

The finding is that weather does not distinguishably improve the forecast at any venue on the
basis the model would actually have. At all three venues the no-weather arm shares the ninety
percent confidence set with every weather arm, including the horizon-matched one, so the
exogenous path is, on this estate and this evidence, decoration. The covariate-quality optimism
the study set out to quantify, the archived-near-valid basis flattering the model relative to
the horizon-matched one, is real and grows with lead, but only at Ellel, where the exogenous
entrant is neither served nor better than no weather at all. At the Beer Hall, the one venue
whose served model is the exogenous one, the archived and the horizon-matched bases are
statistically inseparable, and the horizon-matched arm is if anything nominally the better of
the two. The practical consequence is reassuring and slightly deflating at once: the published
exogenous numbers at the served venue are not inflated by covariate quality, because at that
venue weather barely matters, and perfect knowledge of the forecast period's weather is worth no
more than a seven-day-ahead forecast of it.

\section{A common verdict}
\label{sec:results-common}

The two studies rhyme. In each, an elaboration of the forecaster that is well motivated in the
literature and easy to believe in, cross-series learning in the first and a richer weather
covariate in the second, fails to beat a simpler incumbent once it is measured under a
pre-registered rule with a multiplicity-aware test. Neither result is a failure of the method in
general; each is a statement about this estate, at this scale, on this evidence. The value of
reporting them at full prominence, rather than searching for the venue or the metric on which the
elaboration wins, is that the model confidence set makes the width of a set the finding. Where
the data cannot separate the grouped from the univariate forecaster, or the weather-fed from the
weather-free one, the wider set is the answer, and a spuriously precise winner would say less than
the honest null.
